连续平面波展开已经固定真空光子的两个偏振。在腔和纳米结构中还要回答两个实际问题：一个光子的场幅度如何由能量归一化决定，以及具有同一平均能量的光态为何产生不同交换统计。

\subsubsection{无损正常模的光子归一化}
\label{sec:feqo-quantized-field}
\begin{proposition}[无损正常模的每光子归一化]
在无损、无色散、封闭的非磁性腔中，设实对称正定介电张量为 $\boldsymbol\varepsilon_r(\mathbf r)$。取频率 $\omega_\lambda>0$ 的实驻波模 $\mathbf f_\lambda$ 满足
\begin{align}
 \nabla\times\nabla\times\mathbf f_\lambda
 &=\frac{\omega_\lambda^2}{c^2}\boldsymbol\varepsilon_r\mathbf f_\lambda,
 \qquad \nabla\cdot(\boldsymbol\varepsilon_r\mathbf f_\lambda)=0,\nonumber\\
 \int\dd^3r\,\mathbf f_\lambda\cdot\boldsymbol\varepsilon_r\mathbf f_{\lambda'}
 &=\delta_{\lambda\lambda'}.
 \label{eq:feqo-normal-mode-energy-normalization}
\end{align}
边界取理想导体条件或其他使旋度分部积分表面项为零的自伴条件。单模坐标 $\mathbf A=Q_\lambda\mathbf f_\lambda/\sqrt{\varepsilon_0}$ 给
\begin{equation}
 H_\lambda=\tfrac12(P_\lambda^2+\omega_\lambda^2Q_\lambda^2)
 =\hbar\omega_\lambda(a_\lambda^\dagger a_\lambda+\tfrac12),
 \label{eq:feqo-single-mode-hamiltonian}
\end{equation}
并且该模的正频电场是
\begin{equation}
 \hat{\mathbf E}^{(+)}_\lambda(\mathbf r,t)
 =\ii\sqrt{\frac{\hbar\omega_\lambda}{2\varepsilon_0}}\,
 \mathbf f_\lambda(\mathbf r)a_\lambda\ee^{-\ii\omega_\lambda t}.
 \label{eq:feqo-single-mode-field}
\end{equation}
这些表达式只在上述能量与自伴条件下适用；有色散时需加入材料能量，有损时需保留连续环境。
\end{proposition}

\begin{proof}
令 $\mathbf E=-\dot Q_\lambda\mathbf f_\lambda/\sqrt{\varepsilon_0}$、$\mathbf B=Q_\lambda\nabla\times\mathbf f_\lambda/\sqrt{\varepsilon_0}$。模方程乘 $\mathbf f_\lambda$ 并积分，边界项为零，得到
\begin{equation*}
 \int\dd^3r\,|\nabla\times\mathbf f_\lambda|^2
 =\frac{\omega_\lambda^2}{c^2}
 \int\dd^3r\,\mathbf f_\lambda\cdot\boldsymbol\varepsilon_r\mathbf f_\lambda
 =\frac{\omega_\lambda^2}{c^2}.
\end{equation*}
电磁能 $\frac12\int(\varepsilon_0\mathbf E\cdot\boldsymbol\varepsilon_r\mathbf E+\mu_0^{-1}|\mathbf B|^2)\dd^3r$ 因而为 $(\dot Q_\lambda^2+\omega_\lambda^2Q_\lambda^2)/2$，正则动量为 $P_\lambda=\dot Q_\lambda$。使用\eqnrefs{eq:sho-ladder-r3} 的反解
\begin{equation*}
 Q_\lambda=\sqrt{\frac{\hbar}{2\omega_\lambda}}(a_\lambda+a_\lambda^\dagger),
 \qquad P_\lambda=-\ii\sqrt{\frac{\hbar\omega_\lambda}{2}}(a_\lambda-a_\lambda^\dagger)
\end{equation*}
得到 Hamiltonian；$a_\lambda(t)=a_\lambda\ee^{-\ii\omega_\lambda t}$ 再给电场。由于 $[\mathbf f_\lambda]=\mathrm{m^{-3/2}}$，\eqnrefs{eq:feqo-single-mode-field} 的单位为 $\mathrm{V\,m^{-1}}$。
\end{proof}

\subsubsection{正常模与横向 Green 谱}
\label{sec:feqo-green-spectral}
\begin{proposition}[正常模与横向 Green 谱的对应]
对上述封闭系统，传播子空间中的 Green 核为
\begin{equation}
 \mathbf G_T^R(\mathbf r,\mathbf r';\omega)
 =c^2\sum_\lambda
 \frac{\mathbf f_\lambda(\mathbf r)\otimes\mathbf f_\lambda^*(\mathbf r')}
 {\omega_\lambda^2-(\omega+\ii0^+)^2}.
 \label{eq:feqo-green-mode-expansion}
\end{equation}
在 $\omega>0$，
\begin{equation}
 \operatorname{Im}_{H}\mathbf G_T^R
 =\frac{\pi c^2}{2}\sum_\lambda
 \frac{\mathbf f_\lambda\otimes\mathbf f_\lambda^*}{\omega_\lambda}
 \delta(\omega-\omega_\lambda).
 \label{eq:feqo-green-mode-positive-spectrum}
\end{equation}
量子场与经典 Green 张量使用同一组模式及能量归一化；完整 Green 张量还包含由 Gauss 约束确定的纵向部分。
\end{proposition}

\begin{proof}
在加权内积 $\langle\mathbf u,\mathbf v\rangle_\varepsilon=\int\mathbf u^*\cdot\boldsymbol\varepsilon_r\mathbf v\,\dd^3r$ 下展开任意物理横向场。Maxwell 算符作用于单个模给
\begin{equation*}
 \mathsf L_{\rm EM}(\omega)\mathbf f_\lambda
 =\frac{\omega_\lambda^2-\omega^2}{c^2}\boldsymbol\varepsilon_r\mathbf f_\lambda.
\end{equation*}
逐模求逆并用完备关系得到谱和。$\operatorname{Im}(x-\ii0^+)^{-1}=\pi\delta(x)$ 和正频根的 Jacobian $1/(2\omega_\lambda)$ 给第二式。

把电场模展开中的真空收缩 $\langle0|a_\lambda a_{\lambda'}^\dagger|0\rangle=\delta_{\lambda\lambda'}$ 与谱和比较，可得正频谱分量的真空相关
\begin{equation*}
 \langle0|\hat{\mathbf E}^{(+)}(\mathbf r,\omega)
 \hat{\mathbf E}^{(-)}(\mathbf r',\omega')|0\rangle
 =\frac{\hbar\mu_0\omega^2}{\pi}
 \operatorname{Im}_{H}\mathbf G_T^R(\mathbf r,\mathbf r';\omega)
 \delta(\omega-\omega').
\end{equation*}
离散腔中该式在频率分布意义下读取；它也是后面连续储库涨落公式的无损正常模极限。
\end{proof}

\begin{exbox}[周期盒、连续归一化与每体积的光子态密度]
在边长 $L$、体积 $V=L^3$ 的真空周期盒中，$\boldsymbol\kappa=2\pi\mathbf n/L$，每个非零波矢有两个横向偏振。离散模满足 $[a_{\boldsymbol\kappa\lambda},a_{\boldsymbol\kappa'\lambda'}^\dagger]=\delta_{\boldsymbol\kappa\boldsymbol\kappa'}\delta_{\lambda\lambda'}$。每个格点占据波矢体积 $(2\pi)^3/V$，因此
\begin{equation*}
 \sum_{\boldsymbol\kappa}\longrightarrow\frac{V}{(2\pi)^3}\int\dd^3\kappa,
 \qquad \delta_{\boldsymbol\kappa\boldsymbol\kappa'}
 \longrightarrow\frac{(2\pi)^3}{V}\delta^{(3)}(\boldsymbol\kappa-\boldsymbol\kappa').
\end{equation*}
与\eqnrefs{eq:maxwell-mode-r2} 相同的连续算符是 $a_\lambda(\boldsymbol\kappa)=\sqrt V\,a_{\boldsymbol\kappa\lambda}$，故其对易子为 $(2\pi)^3\delta^{(3)}$，离散场前的 $V^{-1/2}$ 正好转成连续积分系数。若改用没有 $(2\pi)^3$ 的 delta 对易子，必须同时缩放算符和场系数，不能只改积分测度。

真实单光子态由归一化波包构造：
\begin{equation}
 |1_F\rangle=\sum_{\lambda=1}^2\int\frac{\dd^3\kappa}{(2\pi)^3}
 F_\lambda(\boldsymbol\kappa)a_\lambda^\dagger(\boldsymbol\kappa)|0\rangle,
 \qquad \sum_\lambda\int\frac{\dd^3\kappa}{(2\pi)^3}|F_\lambda|^2=1.
 \label{eq:feqo-normalized-photon-wavepacket}
\end{equation}
对易关系给 $\langle1_F|1_F\rangle=1$，平均激发能为
$\sum_\lambda\int\hbar c|\boldsymbol\kappa|\,|F_\lambda|^2\dd^3\kappa/(2\pi)^3$；单个确定波矢的态则是不可归一化的广义本征态。

在频率区间 $[\omega,\omega+\dd\omega]$ 计数两个偏振的格点，得到每单位体积的态密度
\begin{equation}
 \rho_\gamma(\omega)\dd\omega
 =\frac{2}{(2\pi)^3}4\pi\kappa^2\dd\kappa
 =\frac{\omega^2}{\pi^2c^3}\dd\omega.
 \label{eq:feqo-vacuum-photon-density}
\end{equation}
真空 Green 张量的独立核验是
\begin{equation*}
 \frac{2\omega}{\pi c^2}\Tr\operatorname{Im}\mathbf G_0^R(\mathbf r,\mathbf r;\omega)
 =\frac{2\omega}{\pi c^2}\frac{3\omega}{6\pi c}
 =\rho_\gamma(\omega).
\end{equation*}
这里用了\eqnrefs{eq:feqo-vacuum-dipole-radiation} 附近得到的真空局域谱；它把模式计数、每光子能量和 Green 响应的归一化连在一起。
\end{exbox}

\subsubsection{相干态与单模热态的涨落}
\label{sec:feqo-coherent-thermal}
\begin{definition}[相干态与单模热态]
相干态是单位范数的湮灭算符本征态 $a|\alpha\rangle=\alpha|\alpha\rangle$，$\alpha\in\mathbb C$。单模热态是温度 $T>0$ 下的 Gibbs 密度算符 $\rho_{\rm th}=\ee^{-H_\lambda/(k_BT)}/\Tr\ee^{-H_\lambda/(k_BT)}$。二者在数态基中分别为
\begin{align}
 |\alpha\rangle&=\ee^{-|\alpha|^2/2}\sum_{n=0}^\infty\frac{\alpha^n}{\sqrt{n!}}|n\rangle,
 \label{eq:feqo-coherent-state}\\
 \rho_{\rm th}&=(1-x)\sum_{n=0}^\infty x^n|n\rangle\langle n|,
 \qquad x=\ee^{-\hbar\omega_\lambda/(k_BT)},\quad
 \bar n=\frac{x}{1-x}.
 \label{eq:feqo-thermal-state}
\end{align}
相干态的平均振幅非零，热态则 $\langle a\rangle=0$，即使二者有相同 $\bar n$。
\end{definition}

对单模相干态，\eqnrefs{eq:feqo-single-mode-field} 给平均正频场幅度
$\mathbf e_+=\ii\sqrt{\hbar\omega_\lambda/(2\varepsilon_0)}\,\mathbf f_\lambda\alpha$。
完整实场是 $\mathbf e_+\ee^{-\ii\omega_\lambda t}+\mathbf e_+^*\ee^{\ii\omega_\lambda t}$。
若写成 $\operatorname{Re}[\boldsymbol{\mathcal E}\ee^{-\ii\omega_\lambda t}]$，则 $\boldsymbol{\mathcal E}=2\mathbf e_+$；正频幅度与实场复幅度之间的这个二倍关系同时影响电子耦合和时间平均功率。

\begin{proposition}[数涨落、正交分量与二阶相关]
令 $N=a^\dagger a$、$X=(a+a^\dagger)/\sqrt2$、$Y=(a-a^\dagger)/(\ii\sqrt2)$，则 $[X,Y]=\ii$。相干态满足
\begin{equation}
 \langle N\rangle=|\alpha|^2,\qquad
 \operatorname{Var}N=|\alpha|^2,\qquad
 \operatorname{Var}X=\operatorname{Var}Y=\tfrac12.
 \label{eq:feqo-coherent-variances}
\end{equation}
热态满足
\begin{equation}
 \operatorname{Var}N=\bar n(1+\bar n),\qquad
 \operatorname{Var}X=\operatorname{Var}Y=\bar n+\tfrac12.
 \label{eq:feqo-thermal-variances}
\end{equation}
对 $\bar n>0$，零延迟正规序相关 $g^{(2)}(0)=\langle N(N-1)\rangle/\bar n^2$ 依次是相干态的 $1$、热态的 $2$ 和数态 $|n\rangle$ 的 $1-1/n$。
\end{proposition}

\begin{proof}
相干态递推 $c_{n+1}\sqrt{n+1}=\alpha c_n$ 与范数固定 $c_0=\ee^{-|\alpha|^2/2}$，因此数概率是 Poisson 权重。对其产生函数
\begin{equation*}
 F_\alpha(z)=\langle z^N\rangle=\ee^{|\alpha|^2(z-1)}
\end{equation*}
在 $z=1$ 求一、二次导数，得 $\langle N\rangle=|\alpha|^2$ 和 $\langle N(N-1)\rangle=|\alpha|^4$。

热态中零点能在配分函数归一化时相消，故几何权重的产生函数为 $F_T(z)=(1-x)/(1-xz)$。其导数给 $\langle N\rangle=x/(1-x)$、$\langle N(N-1)\rangle=2x^2/(1-x)^2$。由 $N^2=N(N-1)+N$ 得两种方差。

相干态中 $\langle a\rangle=\alpha$、$\langle a^2\rangle=\alpha^2$、$\langle aa^\dagger\rangle=|\alpha|^2+1$。逐项展开 $X^2,Y^2$ 并减去均值平方，剩下各自 $1/2$。因此 $\Delta X\Delta Y=1/2$，或 $\Delta Q\Delta P=\hbar/2$。热态的数态对角性使 $\langle a\rangle=\langle a^2\rangle=0$，而 $\langle aa^\dagger\rangle=\bar n+1$，于是两个正交分量方差均为 $\bar n+1/2$。
\end{proof}

\begin{exbox}[同为四个光子的平均能量，交换统计仍不同]
比较 $|4\rangle$、$\alpha=2\ee^{\ii\theta}$ 的相干态，以及 $x=4/5$ 的热态。它们均有 $\langle N\rangle=4$，但数方差依次为 $0,4,20$，$\langle N(N-1)\rangle$ 依次为 $12,16,32$。

若在受控微扰区间内一光子吸收振幅正比于 $a$，两光子吸收振幅正比于 $a^2$，且其他电子核和失谐相同，那么一光子概率采样 $\langle a^\dagger a\rangle=4$，两光子概率则采样 $\langle a^{\dagger2}a^2\rangle=12,16,32$。所以仅匹配平均光子数无法匹配所有电子边带概率；进入强耦合时必须保留完整联合幺正演化。
\end{exbox}

\begin{definition}[电子电流与量子辐射场的耦合]
在 Coulomb 规范中，横向辐射的线性耦合为
\begin{equation}
 H_{I,T}(t)=-\int\dd^3r\,
 \hat{\mathbf j}_{\rm ch,T}(\mathbf r,t)\cdot\hat{\mathbf A}_T(\mathbf r,t).
 \label{eq:feqo-current-field-coupling}
\end{equation}
Gauss 约束同时产生纵向 Coulomb 能量，不能在具有纵向电荷响应的过程中舍去。单电子电荷流矩阵元是 $q_{\mathrm e}c\Psi_f^\dagger\boldsymbol\alpha\Psi_i$，与\eqnrefs{eq:qed-interaction-lagrangian} 的 QED 顶角使用相同的 Dirac 双线性。
\end{definition}
